% Tandem SRS: LaTeX adaptation of the Meyer (PEGS) requirements structure.
% Reference implementation: Sebastien Mosser, cas-handbook-req-template.
% https://github.com/ace-lectures/cas-handbook-req-template
% Adapted template scaffold: CC BY-NC-SA 4.0.
% Changes: LaTeX implementation, new Tandem authoring prompts, and capstone
% additions identified below. Preserve this path for the editor and build.
\documentclass[12pt]{article}

\usepackage[margin=1in]{geometry}
\usepackage{booktabs}
\usepackage{tabularx}
\usepackage{enumitem}
\usepackage{amsmath,amssymb}

\input{../Comments.text}
\input{../Common.text}

\setlength{\emergencystretch}{3em}
\setcounter{tocdepth}{2}
\setcounter{secnumdepth}{3}

% Replace each authoring prompt with reviewed project content.
\newcommand{\draftnote}[1]{\par\noindent\textit{To complete: #1}\par}
\newcommand{\pending}{\textit{To be completed}}

% Use Meyer's identifiers in headings and cross-references: G.1, E.1, etc.
\newcommand{\requirementsbook}[2]{%
  \clearpage
  \renewcommand{\thesection}{#1}%
  \section{#2}%
  \label{book:#1}%
}

\title{Software Requirements Specification\\\progname\\
  \large Meyer (PEGS) Template}
\author{\authname}
\date{October 18, 2026}

\begin{document}

\maketitle
\thispagestyle{empty}
\clearpage
\pagenumbering{roman}

\section*{Revision History}
\addcontentsline{toc}{section}{Revision History}
\begin{tabularx}{\textwidth}{p{0.17\textwidth}p{0.22\textwidth}X}
\toprule
\textbf{Date} & \textbf{Contributor} & \textbf{Change}\\
\midrule
Oct 18, 2026 & Codex, at Musab Mondal's request &
Replace the blank Volere template with a LaTeX adaptation of Meyer's
structure and add capstone authoring sections.\\
\bottomrule
\end{tabularx}

\section*{About This Document}
\addcontentsline{toc}{section}{About This Document}

This document uses Bertrand Meyer's four-book requirements structure:
Goals (G), Environment (E), System (S), and Project (P). It follows the
structure described in Meyer's
\href{https://se.inf.ethz.ch/requirements/plan.pdf}{Standard Plan for
Requirements} and the
\href{https://github.com/ace-lectures/cas-handbook-req-template}{McMaster
template by S\'ebastien Mosser}. The adapted template scaffold is provided
under \href{https://creativecommons.org/licenses/by-nc-sa/4.0/}{CC BY-NC-SA 4.0}.
The structure has been implemented in LaTeX with new authoring prompts
for Tandem. The source remains at its existing repository path to preserve
the current editor and build setup.

The four books retain the standard chapter structure. The following
capstone additions are included:
\begin{itemize}[noitemsep]
  \item Likely and unlikely changes under Section~\ref{sec:changes}.
  \item Dedicated functional, quality, safety and security requirements,
  plus a critical formal specification, under Section~\ref{sec:functionality}.
  \item A traceability matrix under Section~\ref{sec:traceability} and a
  requirements phase-in plan under Section~\ref{sec:schedule}.
  \item Appendices for auxiliary constants, stakeholder consultation
  evidence, AI use disclosure, and the course reflection questions.
\end{itemize}

\textbf{Draft status.} This file is an authoring framework. Italicized
prompts and table placeholders are to be replaced with reviewed content.
They are not approved requirements or evidence of completed activities.
Use the latest Problem Statement and Goals when completing the scope;
resolve differences with the Development Plan before carrying its earlier
technical approach into this SRS.

Give each requirement a stable identifier, such as \texttt{FR-01},
\texttt{NFR-01}, or \texttt{SR-01}. Record its statement, source, rationale
where needed, and measurable acceptance criterion. Define quantities with
symbolic names in Appendix~\ref{app:constants}. If a section does not
apply, retain it and explain why.

\clearpage
\tableofcontents
\clearpage
\pagenumbering{arabic}

\requirementsbook{G}{Goals}

\subsection{Context and Overall Objective}
\label{sec:objective}
\draftnote{Explain the need Tandem addresses, the intended audience, and
the overall objective. Use the current model-switching scope in the
Problem Statement and Goals.}

\subsection{Current Situation}
\draftnote{Describe the existing workflow and baseline serving behaviour.
Support technical claims with sources or measurements, and distinguish
observations from assumptions.}

\subsection{Expected Benefits}
\label{sec:benefits}
\draftnote{State the benefits expected by stakeholders. Link numbered goals
to the Problem Statement and Goals, and distinguish core goals from
stretch goals. Leave detailed acceptance thresholds in the System book.}

\subsection{Functionality Overview}
\draftnote{Give a short overview of Tandem's capabilities, including context
reuse, ordinary serving and fallback. Keep detailed requirements in
Section~\ref{sec:functionality}.}

\subsection{High-Level Usage Scenarios}
\draftnote{Describe the main interactions from the user's or operator's
viewpoint. Consider a model switch with reusable context, a model switch
without reusable context, and ordinary same-model requests.}

\subsection{Limitations and Exclusions}
\draftnote{Define the product boundary and excluded work. Explain the
rationale for scope decisions and identify stretch capabilities separately.}

\subsection{Stakeholders and Requirements Sources}
\label{sec:stakeholders}
\draftnote{Identify relevant external stakeholder groups and their needs.
Record the documents, consultations and other evidence from which
requirements originate. Link actual consultation records to
Appendix~\ref{app:elicitation}; do not invent meetings or feedback.}

\requirementsbook{E}{Environment}

\subsection{Glossary}
\draftnote{Define all domain terms and acronyms used in this SRS, such as
semantic routing, prefill, decode, KV cache, token and time to first token.}

\subsection{Components}
\label{sec:external-components}
\draftnote{Describe systems and resources outside Tandem's boundary that
interact with it. Include the request source, external routing decisions,
models, tokenizers, serving environment and GPU resources as applicable.
Include a context diagram with labelled inputs and outputs.}

\subsection{Constraints}
\draftnote{State externally imposed limitations and justify each one.
Distinguish environmental constraints from implementation choices recorded
in Section~\ref{sec:technical-choices}.}

\subsection{Assumptions}
\draftnote{Number assumptions about the environment and explain their
rationale. Identify how each assumption will be checked and what changes
if it does not hold. Do not assume away cache or serving failures.}

\subsection{Effects}
\draftnote{Describe how operating Tandem may affect its environment,
including computation, memory use and application response times.}

\subsection{Invariants}
\draftnote{State properties of the environment that must remain true during
operation, with precise conditions and scope. Distinguish these from
internal implementation details.}

\requirementsbook{S}{System}

\subsection{Components}
\label{sec:system-components}
\draftnote{Describe Tandem's major functional responsibilities and their
boundaries. Use only the decomposition needed to organize requirements;
leave detailed architecture and algorithms for the design documents.}

\subsubsection{Likely and Unlikely Changes}
\label{sec:changes}
\draftnote{List and justify likely changes, such as supported model
configurations, and unlikely changes that define stable boundaries.
Assign identifiers so the design can trace to these expectations.}

\subsection{Functionality}
\label{sec:functionality}

\subsubsection{Functional Requirements}
\draftnote{Write individually numbered, observable requirements for request
handling, selected-model behaviour, context identification, cache reuse,
native prefill, failure handling and reporting. Specify both successful
operation and responses to invalid or unsupported inputs.}

\subsubsection{Quality Requirements}
\draftnote{Specify measurable requirements for latency, throughput, output
agreement, resource use, reliability, usability, maintainability and other
applicable qualities. Define the baseline and evaluation conditions.
Use justified symbolic thresholds rather than unsupported numbers.}

\subsubsection{Safety and Security Requirements}
\label{sec:safety}
\draftnote{Include requirements addressing request isolation, cache
integrity, permitted reuse, fallback and other mitigations identified in
the Hazard Analysis. Link each requirement to its hazard or failure case.}

\subsubsection{Formal Specification of Critical Behaviour}
\label{sec:formal}
\draftnote{Select and justify a critical part of Tandem to formalize.
Cache eligibility and fallback are a candidate. Define the domains,
predicates, states, preconditions and required outcomes precisely, then
trace the specification to the corresponding prose requirements.}

\subsection{Interfaces}
\label{sec:interfaces}
\draftnote{Specify every input and output at Tandem's boundary: fields,
types, units, valid ranges, required metadata and error indications.
Include request identifiers, selected models, conversation context,
generation settings, cache metadata, generated text and measurements as
applicable. Describe externally observable interfaces without committing
to unnecessary internal representations.}

\subsection{Detailed Usage Scenarios}
\draftnote{Give each scenario an identifier, actor or triggering event,
preconditions, normal sequence, alternatives, failures and postconditions.
Include accepted reuse, rejected reuse, unavailable context, ordinary
serving and failure recovery. Link scenarios to requirement identifiers.}

\subsection{Prioritization}
\draftnote{Assign and justify requirement priorities. Distinguish the
minimum viable scope, additional requirements and future capabilities.
Use the same identifiers in the phase-in plan.}

\subsection{Verification and Acceptance Criteria}
\label{sec:acceptance}
\draftnote{Define when each requirement is satisfied. Record the workload,
configuration, metric, acceptance threshold and verification method as
appropriate. Cross-reference the V\&V Plan for test procedures.}

\subsubsection{Traceability Matrix}
\label{sec:traceability}
\draftnote{Replace the placeholder row with links from requirements to
goals or stakeholder evidence, hazards where applicable, and acceptance
criteria or planned tests. Ensure all requirements are covered.}

\begin{table}[htbp]
\centering
\small
\caption{Requirements traceability (to be completed)}
\label{tab:traceability}
\begin{tabularx}{\textwidth}{XXXX}
\toprule
\textbf{Requirement ID} & \textbf{Goal or source} &
\textbf{Hazard ID, if applicable} & \textbf{Acceptance criterion or test}\\
\midrule
\pending & \pending & \pending & \pending\\
\bottomrule
\end{tabularx}
\end{table}

\requirementsbook{P}{Project}

\subsection{Roles and Personnel}
\draftnote{Identify responsibilities for requirements, implementation,
review and validation. Reference the Development Plan for team working
arrangements rather than duplicating it.}

\subsection{Imposed Technical Choices}
\label{sec:technical-choices}
\draftnote{Record justified decisions that constrain the solution, such
as an imposed serving engine or supported platform. Distinguish confirmed
constraints from candidates still being evaluated.}

\subsection{Schedule and Milestones}
\label{sec:schedule}
\draftnote{Map requirement identifiers to implementation and validation
milestones, with dates and owners. Reconcile the Development Plan with the
current scope before using its technical milestones.}

\begin{table}[htbp]
\centering
\small
\caption{Requirements phase-in plan (to be completed)}
\label{tab:phase-in}
\begin{tabularx}{\textwidth}{XXXX}
\toprule
\textbf{Milestone and date} & \textbf{Requirement IDs} &
\textbf{Owner} & \textbf{Completion evidence}\\
\midrule
\pending & \pending & \pending & \pending\\
\bottomrule
\end{tabularx}
\end{table}

\subsection{Tasks and Deliverables}
\draftnote{Identify the work products needed for each milestone, including
documentation, feasibility evidence, implementation and evaluation.
Reference the project plan and identify deferred deliverables.}

\subsection{Required Technology Elements}
\draftnote{List required software, models, hardware and supporting tools,
their purpose, relevant compatibility conditions and access limitations.
Record any budget or resource constraints and supporting rationale.}

\subsection{Risk and Mitigation Analysis}
\draftnote{Describe project risks, such as feasibility uncertainty or
hardware access, and the response to each. Reference the separate Hazard
Analysis for system failure modes and trace its mitigation requirements
to Section~\ref{sec:safety}.}

\subsection{Requirements Process and Report}
\draftnote{Describe how requirements are elicited, reviewed, changed and
approved. Link stakeholder evidence, review issues and resulting decisions.
Record unresolved questions and their planned resolution; update this
section as activities are completed.}

\clearpage
\appendix
\renewcommand{\thesection}{\Alph{section}}

\section{Auxiliary Constants}
\label{app:constants}
\draftnote{Define every symbolic threshold used in the SRS, with its value,
units and rationale. Resolve each value before claiming its associated
requirement has a complete acceptance criterion.}

\begin{table}[htbp]
\centering
\small
\caption{Auxiliary constants (to be completed)}
\label{tab:constants}
\begin{tabularx}{\textwidth}{XXXX}
\toprule
\textbf{Symbol} & \textbf{Meaning} & \textbf{Value and units} &
\textbf{Rationale or source}\\
\midrule
\pending & \pending & \pending & \pending\\
\bottomrule
\end{tabularx}
\end{table}

\section{Stakeholder Consultation Evidence}
\label{app:elicitation}
\draftnote{For each actual consultation, record its identifier, date,
participant roles, questions, relevant feedback and resulting requirement
identifiers. Reference meeting records where available. Clearly distinguish
team proposals from stakeholder feedback.}

\clearpage
\section{Generative AI Use Disclosure}

\subsection{AI Tools Used}
OpenAI Codex was used to convert the blank Volere document into a LaTeX
adaptation of the Meyer structure at Musab Mondal's request.
\draftnote{Record the model or version available to the team and any other
AI tools that materially contribute to the completed SRS.}

\subsection{How and Where Used}
Codex reorganized the template into the four Meyer books, wrote authoring
prompts, added capstone traceability and phase-in tables, and retained the
shared team information and reflection questions. It consulted the Meyer
Standard Plan and McMaster reference template. This conversion did not
complete the requirements or supply stakeholder consultation evidence.
\draftnote{Update this disclosure if AI assistance contributes to later
requirements, diagrams, prose or review. Identify the affected sections.}

\subsection{Verification of AI-Generated Content}
\draftnote{Describe the team's actual review of the template adaptation and
any later AI-assisted content. Check technical claims, requirement wording,
formal definitions, references and acceptance criteria. Record compilation
and other checks actually performed; do not present planned review as
completed review.}

The team remains responsible for the accuracy, completeness, consistency
and appropriateness of the submitted document.

\clearpage
\section{Reflection}

\input{../Reflection.text}
\input{../SRS_Reflection.text}

\end{document}
